\section{An orthonormal MPS basis for an approximate low-energy space}

The individual approximation theorem also gives a basis for a space close
to the original inner band. The basis vectors are first approximated with
one common MPS bound. Their actual linear image determines the candidate
space; its dimension and independence are conclusions. Orthogonalizing
inside that same space costs the number of generators in the bond bound.
This is the mathematical basis consequence described in
\cite[line 288 and Theorem~thm:al-ld, lines 1096--1105]{Arad2016RigorousRG}.

\begin{lemma}[The span of a finite indexed MPS family]
    \label{lem:mps_indexed_family_span}
    \lean{TNLean.MPS.AreaLaw.isOpenMPS_of_mem_span_range}
    \leanok
    Let $(\phi_i)_{i\in I}$ be a finite indexed family of open-boundary
    MPS of bond dimension at most $D$. Every vector in their span is
    an MPS of bond dimension at most $|I|D$.
    Repeated or zero family members and empty families are allowed.
\end{lemma}
\begin{proof}\leanok
    Pass to the finite set of values of the family. Its cardinality is
    at most $|I|$. Apply the finite-span cut-rank estimate and enlarge
    the bond bound. No independence hypothesis is needed.
\end{proof}

\begin{theorem}[An actual approximate-band orthonormal MPS basis]
    \label{thm:low_energy_orthonormal_mps_basis}
    \lean{TNLean.MPS.AreaLaw.exists_fixed_bandwidth_low_energy_orthonormal_mps_basis}
    \leanok
    Fix a bandwidth $\eta>0$. There is $K_\eta>0$, chosen before
    the chain and its physical data, with the following property.
    Let $q>0$ and $H=\sum_i h_i$ be the actual positive-contraction
    nearest-neighbour chain. Put
    \[
        G(c)=\mathbf1_{(-\infty,\lambda_0+c]}(H)\mathcal H,\qquad
        T=G(\eta-v),\qquad n=\dim T,\qquad r=\dim G(\eta),
    \]
    where $\lambda_0=\min\sigma(H)$ and $v>0$.
    For $0<\delta\le1$, set
    \begin{align*}
        g        & =\min\{\min(v,\eta/2)/6,1\},                   &
        u        & =1+\log q+\log(1/g),                             \\
        \epsilon & =\frac{\delta}{8(n+1)},                        &
        a        & =16+\log(1/g)+\log\frac{1024(L+1)(n+1)}\delta.
    \end{align*}
    An integer $B\ge1$ is chosen before the target basis and vectors,
    with
    \[
        \log B\le K_\eta
        \left(u^3/g+g^{-1/4}a^{3/4}u\right)(1+\log a).
    \]
    Take the actual standard orthonormal basis $(b_i)_{i=1}^n$ of $T$.
    There is a family of unit MPS $\phi_i$, each of bond dimension at
    most $rB$, with $\|\phi_i-b_i\|^2\le\epsilon/2$.
    For the actual map $F:T\to\mathcal H$ defined by $Fb_i=\phi_i$
    and its range $V$, one has
    \[
        F\text{ injective},\qquad
        V=\operatorname{span}\{\phi_i\},\qquad \dim V=n\le r.
    \]
    The spaces $V$ and $T$ are $\delta$-close in both directions.
    Moreover, $G(\eta)$ is $\delta$-close to $V$: every vector of $V$
    retains at least the fraction $1-\delta$ of its squared norm upon
    projection onto $G(\eta)$.
    Every vector of $V$ is an MPS of bond dimension at most $nrB$.
    In particular, the actual standard orthonormal basis of $V$ is an
    orthonormal MPS family with that bond bound.
\end{theorem}
\begin{proof}\leanok
    \uses{thm:ld_sharp_single_witness,lem:mps_indexed_family_span}
    Apply the proved physical approximation at accuracy $\epsilon$.
    Its common global rank factor is chosen before the target basis
    and works for each of its unit vectors. Since
    $128(L+1)/\epsilon=1024(L+1)(n+1)/\delta$, its accuracy variable
    is exactly the displayed $a$.
    Choose the actual normalized approximant $\phi_i$ to each $b_i$.
    The distance bound retains its complex phase and gives
    \[
        \sum_i\|\phi_i-b_i\|^2
        \le n\epsilon/2\le\delta/16<1.
    \]
    Parseval and finite Cauchy--Schwarz give
    $\|Fx-x\|^2\le(\delta/16)\|x\|^2$ on $T$.
    This proves injectivity and dimension equality, while the basis-defined
    range is exactly the image span. Projection minimality gives forward
    closeness, and equal-dimensional symmetry gives reverse closeness,
    with error at most $\delta/16$, hence with error $\delta$.
    Spectral cutoff inclusion gives $T\subseteq G(\eta)$, which
    proves $n\le r$ and the asserted outer-band direction of closeness.
    Apply Lemma~\ref{lem:mps_indexed_family_span} to the same family:
    every vector of its range, including each vector of its standard
    orthonormal basis, has bond at most $nrB$.
\end{proof}

If $T=0$, the constructed range is zero and both basis index sets are empty;
the argument uses the same finite sums. Scalar local spaces, empty chains
and every positive deficit are included. The number of output basis vectors
is exactly $\dim G(\eta-v)$, not a claim of $r$ low-energy eigenstates.
Under Assumption (LD), $n\le r=\operatorname{poly}(L)$, so the extra
multiplicity factor and logarithm have their stated polynomial and
logarithmic sizes. This is mathematical basis existence. It does not
identify the family with the output of the probabilistic algorithm of
\cite[Theorem~t:alg2, lines 1869--1874]{Arad2016RigorousRG}, and asserts
neither its running time nor its advertised energy-output guarantee.
The basis vectors are not claimed to be energy eigenstates.
